\section{Newtonian three-dimensional bubble}
\label{sec:newtonian}

\status{Uniform mesh is \textbf{closed at both resolutions} --- $25\times25\times50$
and the reference $50\times50\times100$, 1501 steps each, in both methods. Case~2
has not been run. The adaptive half now has a working dynamic remeshing loop
(Section~\ref{sec:amr}) but no production run.}

\subsection{Case definition}

A bubble of radius $R = 0.25$ centred at $(0.5,0.5,0.5)$ in a tank
$[0,1]\times[0,1]\times[0,2]$, no-slip on all walls, zero initial velocity,
pressure and polymeric tensor, integrated to $t = 3$. Two parameter sets,
Table~\ref{tab:params}, taken from the reference work and originally due to
Adelsberger et al.

\begin{table}[htbp]
\centering
\caption{Physical and dimensionless parameters.}
\label{tab:params}
\small
\adjustbox{max width=\linewidth}{%
\pgfplotstabletypeset[
  col sep=space, string type,
  every head row/.style={before row=\toprule, after row=\midrule},
  every last row/.style={after row=\bottomrule},
]{\dat{caso-parametros.dat}}
}
\end{table}

\subsection{Uniform mesh}

\subsubsection{The methods converge toward each other}

This is the central result of the section, and it is a statement about both
methods rather than about either one: refining the mesh brings them together.

\begin{table}[htbp]
\centering
\caption{Difference between the two methods at $t=3$, at both resolutions.}
\label{tab:convergence}
\small
\begin{adjustbox}{max width=\linewidth}
\pgfplotstabletypeset[
  col sep=space, string type,
  columns/grandeza/.style={column name={Quantity}, column type={l}},
  columns/n25/.style={column name={$25\times25\times50$}, column type={r}},
  columns/n50/.style={column name={$50\times50\times100$}, column type={r}},
  columns/fator/.style={column name={Factor}, column type={r}},
  every head row/.style={before row=\toprule, after row=\midrule},
  every last row/.style={after row=\bottomrule},
]{\dat{c1-convergencia.dat}}
\end{adjustbox}
\end{table}

What makes this more than a coincidence is the \emph{direction} each method
moves. The VOF centre of mass falls from $1.4903$ to $1.4766$; the
front-tracking one rises from $1.4615$ to $1.4678$. They approach from opposite
sides toward roughly $1.47$. A single method converging says only that it is
consistent; two independent representations converging from opposite sides
toward the same value is evidence that the value is the right one.

\begin{table}[htbp]
\centering
\caption{Front-tracking against VOF at $t=3$, both resolutions. Shape
quantities from the extracted surfaces by the same routine; mass from the field
integral for VOF and from the surface volume for front-tracking.}
\label{tab:resolution}
\small
\begin{adjustbox}{max width=\linewidth}
\pgfplotstabletypeset[
  col sep=space, string type,
  columns/grandeza/.style={column name={Quantity at $t=3$}, column type={l}},
  columns/vof25/.style={column name={VOF $N{=}25$}, column type={r}},
  columns/ft25/.style={column name={FT $N{=}25$}, column type={r}},
  columns/vof50/.style={column name={VOF $N{=}50$}, column type={r}},
  columns/ft50/.style={column name={FT $N{=}50$}, column type={r}},
  every head row/.style={before row=\toprule, after row=\midrule},
  every last row/.style={after row=\bottomrule},
]{\dat{c1-resolucao.dat}}
\end{adjustbox}
\end{table}

Both methods reproduce the \textbf{ellipsoidal regime} the reference reports for
case~1 --- the height-to-equivalent-diameter ratio is below one in both, that
is, flattened along the direction of rise --- and both reproduce the shape of
the published rise-velocity curve: the velocity grows, peaks near $t = 0.9$, and
decreases toward a steady value.

\begin{figure}[htbp]
\centering
\begin{tikzpicture}
\begin{axis}[width=0.47\linewidth, xlabel={$t$}, ylabel={$z_c$},
             legend pos=north west, grid=both, no markers, legend cell align=left]
  \addplot table[x=t, y=zc] {\dat{bolha3d-c1-N25-vof.dat}}; \addlegendentry{VOF $N{=}25$}
  \addplot table[x=t, y=zc] {\dat{bolha3d-c1-N25-ft.dat}};  \addlegendentry{FT $N{=}25$}
  \addplot table[x=t, y=zc] {\dat{bolha3d-c1-N50-vof.dat}}; \addlegendentry{VOF $N{=}50$}
  \addplot table[x=t, y=zc] {\dat{bolha3d-c1-N50-ft.dat}};  \addlegendentry{FT $N{=}50$}
\end{axis}
\end{tikzpicture}\hfill
\begin{tikzpicture}
\begin{axis}[width=0.47\linewidth, xlabel={$t$}, ylabel={$\Psi$},
             legend pos=south west, grid=both, no markers, legend cell align=left]
  \addplot table[x=t, y=Psi] {\dat{bolha3d-c1-N25-vof.dat}}; \addlegendentry{VOF $N{=}25$}
  \addplot table[x=t, y=Psi] {\dat{bolha3d-c1-N25-ft.dat}};  \addlegendentry{FT $N{=}25$}
  \addplot table[x=t, y=Psi] {\dat{bolha3d-c1-N50-vof.dat}}; \addlegendentry{VOF $N{=}50$}
  \addplot table[x=t, y=Psi] {\dat{bolha3d-c1-N50-ft.dat}};  \addlegendentry{FT $N{=}50$}
\end{axis}
\end{tikzpicture}
\caption{Centre of mass and sphericity at both resolutions.}
\label{fig:resolution}
\end{figure}

\subsubsection{Mass conservation goes the other way}

Mass is the one quantity that does \emph{not} improve for both methods, and the
two resolutions disagree about which method is better --- which is itself the
finding.

At $25\times25\times50$ the front-tracking drift is $-0.096\%$ against
$+0.396\%$ for VOF, a factor of four in favour of front-tracking. At
$50\times50\times100$ it reverses: $-0.851\%$ against $-0.170\%$, a factor of
five the other way.

Reading the time history explains it. Up to step 400 the two front-tracking runs
are nearly identical ($-7.3\times10^{-4}$ and $-6.8\times10^{-4}$ at step 200,
$-2.5\times10^{-3}$ and $-2.8\times10^{-3}$ at step 400). After that they
diverge: the coarse run peaks at $-2.6\times10^{-3}$ and \emph{recovers} to
$-9.6\times10^{-4}$, while the fine run accumulates monotonically. The VOF drift
likewise changes sign between the resolutions. Both coarse-mesh figures are
therefore partly \textbf{cancellation}, not conservation, and neither should be
quoted as a property of its method.

What survives is the fine-mesh statement: at the reference resolution VOF
conserves mass about five times better than front-tracking, and front-tracking's
drift \emph{grows} with refinement.

\subsubsection{The drift that grows is an artefact of the comparison}
\label{sec:deriva-nsub}

The observation that the triangulation ends with the same vertex and triangle
counts at both resolutions is not a puzzle accompanying the drift: it
\emph{is} the drift's explanation. Both production runs were launched with the
same surface subdivision, so the mesh was refined by a factor of two and the
front was not. The ratio $ds/h$ therefore \emph{doubled} between the two runs,
and the pair never was a mesh-refinement comparison --- it varied two things at
once, and the one it was meant to isolate was not the one that moved.

A controlled pair settles it. Running the same case on two meshes with the
surface resolution \emph{following} the grid, so that $ds/h = 0.75$ in both
($N=20$ with three subdivisions, $N=40$ with four), the drift falls:

\adjustbox{max width=\linewidth}{%
\pgfplotstabletypeset[
  string type,
  columns/passo/.style={column name={step}},
  columns/aa0/.style={column name={$A/A_0$}},
  columns/n20/.style={column name={$N=20$}},
  columns/cir20/.style={column name={surg.}},
  columns/n40/.style={column name={$N=40$}},
  columns/cir40/.style={column name={surg.}},
  columns/razao/.style={column name={ratio}},
  every head row/.style={before row=\toprule, after row=\midrule},
  every last row/.style={after row=\bottomrule},
]{\dat{ft-duas-resolucoes.dat}}
}

The ratio stays between $3.2$ and $3.7$ throughout; second order would give
$4$, first order $2$. Over the three points before the first surgery --- pure
advection against pure advection --- it is $3.63$, $3.70$ and $3.61$, a mean of
$3.65$, or an observed order of $\mathbf{1.87}$. \textbf{The drift is
discretisation error and vanishes under refinement.} There is no systematic
component: an interpolation bias would not fall with $h$.

\caveat{At step 500 the ratio drops to $3.21$, and there the fine case has
performed 252 surgery operations against the coarse case's 55 --- the
comparison is no longer advection against advection alone. The clean points are
the early ones, and they are the ones that carry the conclusion.}

Surgery was the named suspect, and a second experiment with everything else
held identical acquits it:

\adjustbox{max width=\linewidth}{%
\pgfplotstabletypeset[
  string type,
  columns/passo/.style={column name={step}},
  columns/comcirurgia/.style={column name={with surgery}},
  columns/semcirurgia/.style={column name={without surgery}},
  columns/razao/.style={column name={worse by}},
  every head row/.style={before row=\toprule, after row=\midrule},
  every last row/.style={after row=\bottomrule},
]{\dat{ft-cirurgia.dat}}
}

\textbf{Surgery does not cause the drift; it reduces it by $42\%$.} Without it
the triangles stretch and the deformed surface is worse represented. The two
runs are \emph{identical up to step 200}, which is where the first surgery
occurs --- the check that the only difference between them is the intended one.

What this leaves standing is narrower and more useful than the original claim:
at the reference resolution VOF conserves mass about five times better, and
front-tracking's drift is a discretisation error of order $1.87$ that the
production pair failed to expose because it held the surface resolution fixed
while refining the mesh.

\subsubsection{The shapes}

The quantities above summarise the shape; the shape itself is a planar section.
Figure~\ref{fig:newtshape} cuts each extracted surface by the vertical plane
through the bubble centre --- the same view the reference uses to compare two
codes --- and overlays the two methods.

The contours are drawn \emph{centred on their own centroid}. Overlaying them at
their absolute heights would hide the difference in shape behind the difference
in rise, and the rise is already in Table~\ref{tab:resolution}.

\begin{figure}[htbp]
\centering
\begin{tikzpicture}
\begin{axis}[width=0.46\linewidth, height=0.42\linewidth,
             xlabel={$x - x_c$}, ylabel={$z - z_c$}, axis equal image,
             title={$25\times25\times50$}, grid=both, no markers,
             legend pos=south east, legend cell align=left, font=\small]
  \addplot+[thick] table[x=x, y=z] {\dat{secoes/newt-N25-vof-c.dat}};
    \addlegendentry{VOF}
  \addplot+[thick, dashed] table[x=x, y=z] {\dat{secoes/newt-N25-ft-c.dat}};
    \addlegendentry{FT}
\end{axis}
\end{tikzpicture}\hfill
\begin{tikzpicture}
\begin{axis}[width=0.46\linewidth, height=0.42\linewidth,
             xlabel={$x - x_c$}, ylabel={$z - z_c$}, axis equal image,
             title={$50\times50\times100$}, grid=both, no markers,
             legend pos=south east, legend cell align=left, font=\small]
  \addplot+[thick] table[x=x, y=z] {\dat{secoes/newt-N50-vof-c.dat}};
    \addlegendentry{VOF}
  \addplot+[thick, dashed] table[x=x, y=z] {\dat{secoes/newt-N50-ft-c.dat}};
    \addlegendentry{FT}
\end{axis}
\end{tikzpicture}
\caption{Newtonian case 1 at $t=3$: vertical section through the bubble centre,
front-tracking against VOF, at both resolutions. Contours centred on their own
centroid so that the comparison is of shape.}
\label{fig:newtshape}
\end{figure}

Both are oblate, as the ellipsoidal regime requires. The visible difference is
not in the overall form but in the \emph{resolution of the contour}: the
marching-cubes section carries 97 points at the coarse mesh against 185 for the
triangulated surface, because the extraction lattice is the fluid mesh while the
triangulation is refined by surgery independently of it. At
$50\times50\times100$ the two carry 188 and 196, and the contours lie close
enough that the faceting of each is the larger visible effect.

\paragraph{The instrument was cleared, not assumed.}
The VOF centre of mass is computed from the marching-cubes surface, which
encloses about 2.3\% less volume than the field at the coarse resolution; that
could bias the centroid. An independent probe computes the centroid directly
from the field and gives $1.490212$ against $1.4903$ from the surface ---
agreement to \num{1e-4}. The volume bias does not translate into a position
bias.

\paragraph{Topological soundness.}
121 output frames per method per resolution, \emph{all} surfaces closed (zero
orphan edges, $\chi = 2$). In the front-tracking runs the triangulation grew from
5120 to 5558 triangles, with 8824 and 9604 surgery operations and a
partition-of-unity residual of \num{2.1e-14}.

\subsection{Case 2: topological change}
\missing{$\rho_1/\rho_2 = 1000$, $\mu_1/\mu_2 = 100$, $\sigma = 1.96$. Lower
surface tension favours break-up, which is where an explicit surface with
surgery and a volume-fraction field should differ most. Not run.}

\subsection{Adaptive mesh}
\missing{blocked on Section~\ref{sec:amr}}
